\documentclass[11pt,a4paper]{article}
\usepackage[utf8]{inputenc}
\usepackage[T1]{fontenc}
\usepackage[margin=2.3cm]{geometry}
\usepackage{graphicx,booktabs,tabularx}
\usepackage[hidelinks]{hyperref}
\setlength{\parindent}{0pt}
\setlength{\parskip}{6pt}
\setlength{\emergencystretch}{2em}
\begin{document}
\thispagestyle{empty}
{\small Universidad Técnica Estatal de Quevedo \quad Ingeniería de Software}\par
\vspace{12mm}
{\LARGE\bfseries TurnoJusto}\par
{\large Entrega de la práctica AP-U1 · Grupo 2}\par
\medskip
Melanie Michelle Muñoz Vera\\
Aplicaciones Distribuidas (ISR-701)\\
23 de noviembre de 2026
\bigskip

\section*{Repositorio de entrega}
\url{https://github.com/Melanirr/AP-U1-Munoz-Melanie}

El repositorio reúne el prototipo Java, las instrucciones para ejecutar el coordinador y las tres ventanillas, las capturas de consola y los registros de cada proceso.

\begin{tabularx}{\linewidth}{@{}lX@{}}
\toprule
Material & Ubicación en el repositorio\\ \midrule
Guía de ejecución & \texttt{README.md}\\
Informe y fuente & \texttt{docs/informe.pdf} y \texttt{docs/informe.tex}\\
Evidencias & \texttt{docs/img/} y \texttt{logs/}\\
Orden de eventos & \texttt{logs/orden\_total.log}\\
\bottomrule
\end{tabularx}

\section*{Resultado principal}
Con cupo 10, la reproducción normal asignó los números 1 a 10 y rechazó 20 solicitudes. En modo inseguro se aceptaron las 30 solicitudes y todas devolvieron 1. El informe distingue las capturas del laboratorio de la reproducción local de E1 a E6. Queda pendiente la captura inicial de los adaptadores de red de E0.
\end{document}
